\documentclass[10pt,a4paper]{amsart}
\usepackage[margin=19mm]{geometry}
\usepackage{amsmath,amssymb,mathtools}
\usepackage[scr=rsfs,cal=euler]{mathalfa}
\usepackage{xfrac}
\usepackage[unicode,hidelinks,bookmarks=false]{hyperref}
\newcommand{\ie}{i.e.\ }
\newcommand{\cf}{cf.\ }
\newcommand{\resp}{respectively\ }
\newcommand{\Subsection}{\subsection}
\providecommand{\nobreakdash}{\nobreak\hbox{-}}
\setlength{\emergencystretch}{2em}
\multlinegap0pt
\usepackage{mathtools}
\newtheorem{theorem}{Theorem}
\newtheorem{lemma}{Lemma}
\newcommand{\BP}{\mathcal{BP}}
\newcommand{\Gr}{\mathrm{Gr}}
\newcommand{\R}{\mathbb{R}}
\newcommand{\Sph}{\mathbb{S}}
\newcommand{\dd}{\,d}
\title{Extensions of the Busemann-Petty Problem for Arbitrary Measures}
\author{Daniel Galicer, Juli\'an Haddad, Joaqu\'in Singer}
\date{Source: arXiv:2512.13555v3 (CC BY 4.0)}
\begin{document}
\maketitle

\noindent\emph{Source and licence.} Extensions of the Busemann-Petty Problem for Arbitrary Measures. Version arXiv:2512.13555v3 (CC BY 4.0), distributed by arXiv under the Creative Commons Attribution 4.0 International licence, http://creativecommons.org/licenses/by/4.0/, as recorded in the arXiv licence metadata for this exact version and shown on the abstract page https://arxiv.org/abs/2512.13555v3. Full licence notice: \texttt{licenses/arxiv-2512.13555-CC-BY-4.0.txt}.\par\smallskip

\noindent\textit{Editorial scope.} Counted unit: the article's theorem thm:main (source0.tex lines 244--294). The article's complete original proof of it is delivered with it (source0.tex lines 634--729, one \texttt{proof} environment of 96 lines, which the article places away from the statement). No mathematics is written, repaired or re-derived: the statement and the proof are the article's own lines, in the article's own notation, and no retained line is substituted or re-flowed.  Retained uncounted as the source-local context the proof needs: lines 206--209; lines 576--592.  Nothing was omitted from any retained span, so the package carries no omission record.  No retained line contains a citation, so the wrapper carries no bibliography and no bibliography entry.  No retained line contains a figure environment, an image-inclusion command or drawing code, so no image is delivered and the package declares no asset dependency.  Editorial formatting only, in addition: an AMS \texttt{amsart} preamble that restates the article's own declarations and macros used by the retained lines, namely the environments \texttt{theorem}, \texttt{lemma} on separate counters, together with the standard packages those lines need.  The retained environments print in this document as Theorem 1, Lemma 1 in order of appearance, whereas the article prints the same items with the numbering of its own full document.  The complete native source of the article as distributed in the arXiv e-print for this exact version is bundled unchanged as \texttt{sources/191040/source0.tex}.
\begin{equation}\label{eq:BP-definition}
 \int_{\Sph^{n-1}}\rho_D(\theta)^k \varrho(\theta)\dd\theta
 =\int_{\Gr_{n-k}}\mathcal R_{n-k}\varrho(H)\dd\eta_D(H)
\end{equation}

\begin{lemma}\label{lem:continuity}
Let $f,g:\R^n\to\R$ be continuous and let $K,L$ be star bodies. Define
$Q:\Sph^{n-1}\to\R$ by
\[
 Q(\theta)=\int_{\rho_K(\theta)}^{\rho_L(\theta)}w(r\theta)\dd r,
\]
where
\[
 w(x)=
 \begin{cases}
  f(x)&\text{if }x\in K\setminus L,\\
  g(x)&\text{if }x\in L\setminus K,\\
  0&\text{otherwise.}
 \end{cases}
\]
Then $Q$ is continuous.
\end{lemma}

\begin{theorem}\label{thm:main}
Let $1\leq k<n$, let $K,L\subset\R^n$ be origin-symmetric star bodies, and
let $\mu_n,\mu_{n-k},\nu_n,\nu_{n-k}$ be four measures with non-negative even densities
$f_n,f_{n-k},g_n,g_{n-k}$, respectively. Assume that $f_{n-k}$ and
$g_{n-k}$ are continuous and that $f_n,g_n$ are integrable on
$K\mathbin\triangle L$.

On $K\mathbin\triangle L$, set
\[
 u_j(x)=
 \begin{cases}
  f_j(x),&x\in K\setminus L,\\
  g_j(x),&x\in L\setminus K,
 \end{cases}
 \qquad j\in\{n-k,n\}.
\]
Suppose that there are finite, even, measurable functions
$a,b:\R^n\setminus\{0\}\to\R$ such that $a(r\theta)$ is non-decreasing and
$b(r\theta)$ is non-increasing in $r>0$, for every
$\theta\in\Sph^{n-1}$, and
\begin{equation}\label{eq:factorization}
 u_{n-k}(x)=|x|^k u_n(x)\bigl(a(x)+b(x)\bigr)
 \quad\text{for almost every }x\in K\mathbin\triangle L.
\end{equation}
Assume further that
\begin{equation}\label{eq:M}
 c(\theta)
 =a(\rho_L(\theta)\theta)+b(\rho_K(\theta)\theta)
\end{equation}
is positive and continuous and let $M$ be the auxiliary star body with
\[
 \rho_M(\theta)=c(\theta)^{-1/k}.
\]
If
\begin{equation}\label{eq:section-hypothesis}
 \mu_{n-k}\bigl((K\setminus L)\cap H\bigr)
 \leq
 \nu_{n-k}\bigl((L\setminus K)\cap H\bigr)
 \qquad\text{for every }H\in\Gr_{n-k},
\end{equation}
then
\begin{equation}\label{eq:main-conclusion}
 \mu_n(K\setminus L)
 \leq d_{\mathrm{BM}}(M,\BP_k^n)^k\,
       \nu_n(L\setminus K).
\end{equation}
In particular, if $M\in\BP_k^n$, then
\[
 \mu_n(K\setminus L)\leq\nu_n(L\setminus K).
\]
\end{theorem}

\begin{proof}[Proof of Theorem~\ref{thm:main}]
Define
\begin{equation}\label{eq:qtheta}
 q(\theta)=
 \int_{\rho_K(\theta)}^{\rho_L(\theta)}
 r^{n-k-1}u_{n-k}(r\theta)\dd r.
\end{equation}
By Lemma~\ref{lem:continuity}, applied to
$|x|^{n-k-1}f_{n-k}(x)$ and
$|x|^{n-k-1}g_{n-k}(x)$, the function $q$ is continuous and even. The
hypothesis
\[
 \mu_{n-k}\bigl((K\setminus L)\cap H\bigr)
 \leq
 \nu_{n-k}\bigl((L\setminus K)\cap H\bigr)
 \qquad\text{for every }H\in\Gr_{n-k}
\]
is equivalent to
\begin{equation}\label{eq:section-hypothesis2}
 0\leq
 \int_{(L\setminus K)\cap H}u_{n-k}(x)\dd x
 -\int_{(K\setminus L)\cap H}u_{n-k}(x)\dd x
 \qquad\text{for every }H\in\Gr_{n-k}.
\end{equation}
Using polar coordinates in each $H\in\Gr_{n-k}$ shows that
\eqref{eq:section-hypothesis2} is equivalent to
\[
 0\leq\mathcal R_{n-k}q(H)
 \qquad\text{for every }H\in\Gr_{n-k}.
\]

Let $D\in\BP_k^n$ and denote by $\eta_D$ its representing
measure. Integrating the last inequality and using \eqref{eq:BP-definition},
we obtain
\begin{equation}\label{eq:positive-D}
 0\leq\int_{\Gr_{n-k}}\mathcal R_{n-k}q(H)\dd\eta_D(H)
 =\int_{\Sph^{n-1}}\rho_D(\theta)^k q(\theta)\dd\theta.
\end{equation}
By \eqref{eq:factorization},
\[
 q(\theta)
 =\int_{\rho_K(\theta)}^{\rho_L(\theta)}
 r^{n-1}u_n(r\theta)\bigl(a(r\theta)+b(r\theta)\bigr)\dd r.
\]
Radial monotonicity gives, with $c$ as in \eqref{eq:M},
\begin{equation}\label{eq:pointwise-bound}
 \int_{\rho_K(\theta)}^{\rho_L(\theta)}
 r^{n-1}u_n(r\theta)\bigl(a(r\theta)+b(r\theta)\bigr)\dd r
 \leq c(\theta)
 \int_{\rho_K(\theta)}^{\rho_L(\theta)}
 r^{n-1}u_n(r\theta)\dd r.
\end{equation}
Indeed, if $\rho_K(\theta)\leq\rho_L(\theta)$, then for every
$r\in[\rho_K(\theta),\rho_L(\theta)]$ we have
$a(r\theta)\leq a(\rho_L(\theta)\theta)$ and
$b(r\theta)\leq b(\rho_K(\theta)\theta)$, so the factor $a+b$ is bounded above
by $c(\theta)$ on the interval of integration, whose orientation is positive.
If $\rho_L(\theta)\leq\rho_K(\theta)$, the same two monotonicity properties
bound $a+b$ from below by $c(\theta)$ on $[\rho_L(\theta),\rho_K(\theta)]$,
and reversing the orientation of the integral yields exactly the same
inequality.

Combining \eqref{eq:positive-D}, \eqref{eq:pointwise-bound}, and
$c=\rho_M^{-k}$ yields
\[
 0
 \leq
 \int_{\Sph^{n-1}}
 \left(\frac{\rho_D(\theta)}{\rho_M(\theta)}\right)^k
 \int_{\rho_K(\theta)}^{\rho_L(\theta)}
 r^{n-1}u_n(r\theta)\dd r\dd\theta,
\]
which, written in Cartesian coordinates, is
\begin{equation}\label{eq:before-BM}
 \int_{K\setminus L}
 \left(\frac{\rho_D(x/|x|)}{\rho_M(x/|x|)}\right)^k
 f_n(x)\dd x
 \leq
 \int_{L\setminus K}
 \left(\frac{\rho_D(x/|x|)}{\rho_M(x/|x|)}\right)^k
 g_n(x)\dd x.
\end{equation}

Now suppose that $\lambda^{-1}D\subseteq M\subseteq D$. Then
\[
 1\leq\frac{\rho_D(\theta)}{\rho_M(\theta)}\leq\lambda
 \qquad(\theta\in\Sph^{n-1}),
\]
and consequently
\[
 \mu_n(K\setminus L)
 \leq\lambda^k\nu_n(L\setminus K).
\]
Taking the infimum over all admissible $D$ and $\lambda$ proves
\eqref{eq:main-conclusion}.
\end{proof}

\end{document}
